% ============================================================================
% CH — THE MERGE: fUfI-ProtoMF                          [slug: fufi-merge]
% STATUS: PLANNED — outline-level only. The merge candidate still passes its
% own design-candidate + scrutiny cycle before anything lands here.
% Lineage steps 6-7. Mirrors the ProtoMF paper's own build order
% (item-side -> user-side -> combined).
% DELTA CHAPTER, SECOND ORDER (Matteo, 2026-07-14): both sides are fully
% built by now — only merge-specific phenomena belong here (double tie under
% grounded factors, which-side-is-free intercept, refund test, cold users).
% Deliberately the THINNEST lineage chapter, and deliberately the least
% detailed skeleton: no candidate planned out yet and zero results — this
% fills in once we get our hands on it. Do not pad it to look like its
% siblings.
% HEADERS REWORKED 2026-07-14 (granularity): 4 -> 3 sections; the build-order
% mirror is now the CHAPTER INTRO (bonbon), not a section.
% Material map: ../outline.md#fufi-merge
% ============================================================================
\cleardoublepage
\chapter{The Merge: (a)fUfI-ProtoMF (6--8)}   % page goal — scaffolding, DELETE before submission
\label{ch:fufi}
\label{sec:fufi-build-order}   % old opener section parked here (chapter intro)
% The "(a)" in the title (Matteo, 2026-07-14): afU very likely lands before
% the merge — the merged model may already carry user-attribute rows.

% CHAPTER INTRO (bonbon, unnumbered): the thesis extends ProtoMF the way
% ProtoMF built itself — I + U -> UI.

\section{The Tie}
\label{sec:fufi-design}
\vault{2026-07-12_1815_ui-merge-intercept-feature-attributable-check-which-side-is-free, 2026-03-18_1234_double-tie-explainability-tradeoff}
% region vault: 2026-07-12_1815_ui-merge-intercept-feature-attributable-check-which-side-is-free -- auto-copy, edit the vault entry, not this block
%| # The intercept on the UI host: present, one live + one inert — and in fUfI it becomes feature-attributable ("check which side is free")
%|
%| *(dc05 review sitting, follow-up to [[2026-07-12_1808_bt-implicit-item-intercept-shift-asymmetry-directives]].
%| User question: "Is this effect present in UI-ProtoMF / is it a risk in any
%| merged candidate we might come up with?" Answer worked through and captured;
%| the fUfI-facing parts also live in the scratchpad merge-baseline entry.)*
%|
%| ## Present in UI-ProtoMF — one live, one inert
%|
%| UI-score = u*ᵀt̂ + ûᵀt* (A1 Eq. 12). Each shifted vector contributes its
%| constant: the **U-half's 1ᵀt̂ is a per-ITEM scalar → live** (rank-active; the
%| model reaches it as (Wᵗᵀ1)·t, a learned linear functional of the item
%| embedding — it controls both the dial and the direction it reads from); the
%| **I-half's 1ᵀû is a per-USER scalar → rank-inert** for item ranking. So the
%| published UI model already carries exactly one popularity-capable implicit
%| intercept, through the half whose explanation machinery never separated it.
%| No flag removes it (it is structural, not the `use_bias` scalar).
%|
%| ## The fUfI twist: the live intercept decomposes over feature words
%|
%| In the merge the projected item is feature-composed, t̂ = Wᵗ(Σ_f e_f + e_ID),
%| so the live intercept splits **exactly**:
%|
%| > 1ᵀt̂ = Σ_f (Wᵗᵀ1)·e_f + (Wᵗᵀ1)·e_ID
%|
%| — per-feature-word bias contributions plus an item-specific residual. Two
%| readings worth keeping verbatim: (a) this is **LightFM's composed bias
%| b_i = Σ_f b_f re-emerging implicitly** — the very channel dc01's deviation 4
%| dropped in its explicit form; (b) it is one concrete realization of the
%| length-refund prediction that "popularity can live in ID-row length through
%| the linear half" (the component along Wᵗᵀ1). Renderer consequence for the
%| merge stage: the baseline line becomes zoomable ("item baseline +1.9: of
%| which department-word +1.2, colour-word +0.3, item-specific +0.4") — class-
%| level popularity gets *named per feature word*, only the residual stays
%| opaque. The channel converts from pure risk into attributable content.
%|
%| ## Risk verdict + the general rule
%|
%| Any merged candidate on the UI host inherits the live channel by
%| construction. By the merge stage we will have: the measured magnitude on both
%| single-sided hosts (dc05's primary B(t) instruments), the renderer's baseline
%| line as standing practice, and the per-word attribution above. **General
%| diagnostic rule for every future candidate's hidden-effects pre-audit:
%| wherever a shifted-cosine side is scored against free per-item coefficients,
%| a live item intercept exists — check which side is free.**
%|
%| [[phase-4]] [[explanations]] [[decisions]] [[experiments]]
% endregion vault: 2026-07-12_1815_ui-merge-intercept-feature-attributable-check-which-side-is-free
% region vault: 2026-03-18_1234_double-tie-explainability-tradeoff -- auto-copy, edit the vault entry, not this block
%| # Double-tie architecture creates compromised embeddings
%|
%| Deep analysis of the `prototypes_double_tie` (user_item_proto) architecture revealed a fundamental explainability tradeoff. Weight sharing forces each side's embedding to serve two conflicting optimization pressures: (1) being well-positioned for cosine similarity with its own prototypes, and (2) containing information that survives linear projection into a useful signal for the other side's prototype space. This compromises the clean one-dimension-one-prototype correspondence that makes single-prototype variants (`item_proto`, `user_proto`) fully explainable. The projected embedding branches are opaque linear combinations, and even the prototype similarity values are computed from embeddings shaped by both pressures. This is a key consideration for Phase 4 feature-aware design — extending the architecture should preserve explainability where possible.
%|
%| [[paper-understanding]] [[decisions]] [[phase-2]]
% endregion vault: 2026-03-18_1234_double-tie-explainability-tradeoff
\vault{2026-08-15_1721_merge-norm-entanglement-shared-rows-two-sided-refund-2x2, 2026-08-15_1651_spectral-scale-gauge-inversion-direction-diversity-metric, 2026-08-17_1823_correction-double-tie-projection-basis-and-ui-effective-count, 2026-08-17_1320_popularity-stance-2x2-figure-plan-base-settings-gate, 2026-08-17_1937_base-settings-decision-shifted-bias-off-cosbias2x2-results}
% region vault: 2026-08-15_1721_merge-norm-entanglement-shared-rows-two-sided-refund-2x2 -- auto-copy, edit the vault entry, not this block
%| # Merge norm entanglement: shared rows weld the popularity channel to the steering — two-sided refund upgrade, host-tie separability contrast, the built-in 2×2, and the expressivity resolution
%|
%| *(Verbatim-capture entry from the F5 reading session, 2026-08-15 — a five-step build-up worked through interactively at Matteo's questions, all specifically about the **fUfI merge stage**. Destined for the hidden-effects section and the merge chapter (`fufi-merge`). Everything here is hypothesis/mechanism derivation — no measurements; designed tests belong to the merge cycle. Siblings: [[2026-08-15_1639_pm-purity-port-groundtruth-target-discussion]], [[2026-08-15_1651_spectral-scale-gauge-inversion-direction-diversity-metric]].)*
%|
%| ## 0. Matteo's opening question and the factual correction
%|
%| Question: at the merge, prototype latents are constructed from shared attribute embeddings, so a prototype's norm is "subject to the universe, not to itself," while the linear part of a side is more free — could that entanglement cause distorting or performance-inhibiting effects?
%|
%| Correction that reframes it: **in our lineage the prototypes are free** `nn.Parameter`s at every stage (grounding is read off via cos(e_f, p_k), unenforced; constructed prototypes were dc02/dc04, both shelved; even under a centroid-*init* scheme they are free after step one). What is "subject to the universe" is the **composed entity embedding** (t_i = Σ e_v + id_i; q_u = mean of basket rows + id_u). And at the double tie "the linear part" is not a separate parameter at all — the tie reuses the same composed embedding, projected through the other side's prototype matrix. So the concern is real, but its locus is the **shared attribute rows**, not the prototypes. (Prototypes stay clean even at the merge: their projection job is norm-sensitive but their cosine job is norm-blind, so the norm serves the one job without compromise.)
%|
%| ## 1. The mechanism: at the merge, each shared row's norm holds two jobs with different masters
%|
%| At fUfI, composed embedding norms become behaviorally meaningful (they feed the linear halves) *and* they are class-tied. A shared row's norm is one scalar serving:
%|
%| - **cosine branch:** relative row norms = mixture weights of every carrying item's direction (which attributes dominate its identity — the interpretable structure);
%| - **linear branch:** row norm = score magnitude class-wide (e.g. "denim items get a boost" — a class-level popularity prior; LightFM's feature-composed bias mechanism arriving uninvited).
%|
%| The optimizer must compromise. Concrete distortion risk: **popularity pressure on the linear branch inflates certain rows' norms, and the inflated norms drag the cosine-branch directions of every carrying item toward the popular attributes** — popularity contaminating the interpretable direction structure through the norm channel, with no bias term anywhere in sight. ("Removing the bias term didn't remove the bias," one model later, wearing a new coat.)
%|
%| **Two-sided upgrade of the length-channel refund hypothesis** ([[2026-07-12_1224_length-channel-refund-hypothesis-hidden-effects-section]]): the registered refund (linear halves make norm/length expressive again at the merge) and this contamination are **the two signs of one mechanism** — the reopened norm channel can refund the angular tax *or* smuggle popularity into the directions, and which sign dominates is not decidable from the armchair. Designed-test material for the merge cycle, with knobs ready if the bad sign shows: explicit feature-composed bias term (give popularity a dedicated home — the dc01 bias-gap fix), ID-row division of labor, centering.
%|
%| ## 2. Host double tie: the same two-jobs structure, but SEPARABLE — the problem is tie × composition, not the tie
%|
%| Host UI-ProtoMF has the identical two-jobs anatomy (score = u*·(tP^u) + (uP^t)·t*; every entity embedding: cosine direction + linear vector; every prototype: cosine anchor + projection axis). But the host is safe because embeddings are per-entity free:
%|
%| - inflating ‖t_i‖ boosts item i's linear half, leaves i's cosine direction untouched (cosine is norm-blind), and moves **no other item** — no shared rows;
%| - popularity lands as clean **per-item intercepts** (the already-catalogued implicit item-intercept channel, [[2026-07-12_1808_bt-implicit-item-intercept-shift-asymmetry-directives]]); direction structure and norm structure optimize quasi-independently.
%|
%| The host's genuine cross-job tension is on **directions** (cosine anchor vs projection axis) — and per the collapse entry's hypothesis that tension is plausibly a *benefit* (redundancy costs projection rank accuracy sees → anti-collapse pressure; coheres with F-DC01-07). Same structural feature, opposite sign, because it acts on directions of free vectors, not norms of shared ones.
%|
%| |  | free embeddings | composed embeddings |
%| |---|---|---|
%| | **no tie** (single-sided) | norm gauge on proto side, free intercept on linear side — nothing to entangle | norms gauge → norm channel **closed** (the tax); angular effects only |
%| | **double tie** | host UI: channel open, per-entity → **separable, benign** | fUfI: channel open, class-tied, doubly employed → **non-separable** |
%|
%| **The built-in 2×2 (measurement gift):** the four cells are four models the lineage has or will have — host U/I (neither), host UI (tie only), fI/fU (composition only), fUfI (both). Host UI is the free control arm: if popularity contamination of direction structure appears at fUfI but not in host UI (identical tie, free embeddings), the tie×composition interaction is isolated with no new architecture. The refund-vs-contamination test ships with its own control group because of how the lineage was staged.
%|
%| ## 3. What "class" means, precisely (Matteo's hammer-down)
%|
%| The channel's primitive unit is one **attribute value** — one row = one value v; its class = the carrier set of v (all items with colour_group=Black). Hence:
%|
%| - **classes overlap, not partition** — each H&M item sits in exactly 5 value-classes (+ its ID singleton); its linear-half exposure = the sum of its five value effects + ID effect;
%| - **conjunctions are not addressable** — no row for black∩denim; boosting the intersection necessarily boosts all black and all denim. The channel's vocabulary = additive combinations of single-value effects (the composition-vs-FM boundary again: no interaction terms, for better and worse);
%| - **cross-item coupling** = norm-weighted signature overlap: share nothing → independent; share all 5 → **signature twins, maximal coupling** (identical linear halves up to ID rows). The entanglement structure is a weighted overlap graph; twins are its cliques; the finest distinction the shared channel can draw is the full 5-value signature; below that, only the ID row differentiates. (The ID/twins braid's "feature rows carry only class-level signal," now with "class" made precise; F-DC01-06 continuum = the coupling graph's clique structure.)
%| - **fU footnote:** user-side membership is graded — exposure to e_v ∝ v's share of the basket (mean-normalized weights); classes are soft. **ml-1m footnote:** variable tag bags (3–72 tokens, raw-bag policy) → tag-rich items sit in more classes, coupling them to more of the catalog (token-mass coupling inside the entanglement graph).
%|
%| ## 4. The expressivity resolution (Matteo: "is norm a less powerful channel — like neurons sharing weights?")
%|
%| Split by arm — the analogy is right in one and instructively wrong in the other:
%|
%| - **noid arm: yes, strictly less expressive.** Linear halves confined to additive spans of V shared rows (~V·d DoF vs the host's N·d); twins forced identical; conjunctions unaddressable. Hard weight sharing in the CNN sense — one filter reused across positions.
%| - **ids arm: NO — the function class is identical to the host's.** t_i = Σe_v + id_i is a **reparameterization**, not a restriction (any target reachable via id_i = target − Σe_v). What changes:
%|   1. **Repricing (partial pooling, the better analogy).** Under weight decay, "boost these 500 denim items" costs the host 500 per-item norms and fUfI **one row**. Class-shaped signals become cheap, idiosyncratic ones cost full price through the ID row — a multilevel model: shared rows = group effects, ID rows = individual effects, weight decay = the pooling prior. For class-aligned popularity the channel is *more* powerful per parameter — simultaneously the cold-start feature ("denim sells well" survives for a new item) and the contamination risk. One repricing, two faces.
%|   2. **Orthogonality loss — the sharpest sentence of the thread.** The host owns a per-item **direction-preserving dilation** (scale t_i: norm moves, direction untouched — what made its popularity channel clean). In fUfI no such move exists: only components can be scaled — scaling a shared row re-steers every carrying item; scaling the ID row rotates the item toward its ID (eroding feature-explained fraction). For any item sharing even one value with anyone, **no parameter move changes norm without changing direction.**
%|
%| **The corrected one-liner:** *noid: strictly less expressive (hard sharing). ids: equally expressive, differently priced — class-level cheap (partial pooling), idiosyncratic full price — and no longer separable: every norm move is also a direction move, because the host's per-entity direction-preserving dilation does not exist in the composed parameterization.* This *derives* the separability collapse from the parameterization rather than asserting it, and plugs into F-DC01-08 (L2-as-canonicalizer), the ids/noid division-of-labor entries ([[2026-07-12_1830_m5-noise-for-crowding-ids-division-of-labor-centering-option]]), and the feature-composed-bias idea as the friendly face of the same repricing.
%|
%| ## 5. Routing
%|
%| - **Hidden-effects inventory** (`hidden-effects`): joins as the merge-stage sibling of the braid — "norm channel welded to the steering"; the refund hypothesis is now explicitly two-sided.
%| - **Merge chapter** (`fufi-merge`): the 2×2 as the experiment design skeleton for the refund/contamination test; host-dependence of the effect doubles as exposition structure (teach the host's clean channel, then show composition welding it shut).
%| - **Instrument tie-in:** the direction-contamination arm is measured with the direction-diversity/purity instruments of the two sibling entries (per-field P_m profile, normalized-spectrum measures) — popularity axis vs attribute axes on the same PCs (R²-style) is the natural adjudicator.
%| - Scrutiny-evidence-not-thesis applies; nothing here is scheduled before the merge cycle's own DC + scrutiny steps.
%|
%| [[explanations]] [[evaluation]] [[decisions]] [[phase-4]]
% endregion vault: 2026-08-15_1721_merge-norm-entanglement-shared-rows-two-sided-refund-2x2
% region vault: 2026-08-15_1651_spectral-scale-gauge-inversion-direction-diversity-metric -- auto-copy, edit the vault entry, not this block
%| # F5's spectral-scale metric doesn't port — instrument invariances must match readout invariances; norm becomes meaningful again at the double tie
%|
%| *(Verbatim-capture entry from the F5 reading session, sibling of [[2026-08-15_1639_pm-purity-port-groundtruth-target-discussion]]. Status per Matteo: **awareness-level, not a commitment** — we are not 100% sure we will use this metric or something like it, but "we possibly need a metric capturing direction diversity at that level" is now a recorded need. The double-tie norm observation below is flagged by Matteo as extremely interesting.)*
%|
%| ## The question
%|
%| F5 introduces σ_geo (geometric mean of the singular values of the centered embedding matrix) to measure whether a representation has expanded beyond a compressed init ball, explicitly *criticizing effective rank for being scale-invariant* — a tiny isotropic noise ball can look structurally rich to effective rank while being too small to move any logit. Matteo asked: is this relevant in our space at all? We only read out cosine similarities — every prototype could be normalized onto the unit sphere; wouldn't "spread around the ball" be the right measurement instead?
%|
%| ## Answer: correct — and F5's own critique inverts in our space
%|
%| **The gauge argument (confirmed, not just theoretical).** Our prototypes are read through shifted cosine and nothing else: scaling a prototype changes no output. Norm is pure gauge; the only force on it is weight decay — and the S0.7 collapse measurements showed exactly that (all prototype norms decayed to near-identical values, 0.123–0.135, cv 1.6%; see [[2026-07-15_1112_prototype-collapse-is-a-self-built-popularity-bias]]). Any norm-sensitive metric on our prototypes measures optimizer dynamics, not model geometry.
%|
%| **The inversion.** F5's effective-rank critique is *space-relative*, not absolute:
%| - In THEIR space (raw inner-product readout ⟨h, x⟩) absolute scale is behaviorally meaningful → scale-invariant metrics mislead → σ_geo is the fix.
%| - In OUR space (cosine readout) scale is gauge → scale-*sensitive* metrics (σ_geo) mislead — they would report weight-decay behavior — → scale-invariant spectral measures are exactly right.
%|
%| **General lesson (the transferable part): a geometry instrument must match the invariances of the model's readout.** Quoting F5's σ_geo argument in our setting would import a correction for a disease we don't have; conversely their critique of effective rank must never be repeated against our normalized-spectrum measurements.
%|
%| ## The direction-diversity instrument menu (candidates, not commitments)
%|
%| Row-normalize the prototype matrix first (the gauge fix), then:
%| 1. **Effective direction count** — pairwise cosine with duplicate threshold (cos > 0.999). Already the committed SC.8 instrument from the collapse finding; bluntest, most legible ("76 nominal, ~2 effective").
%| 2. **Effective rank / spectral entropy of the normalized matrix** — the smooth, threshold-free version (5 tight clusters → spectrum ≈ rank 5). After normalization the Frobenius norm is pinned at √K, so scale-blindness is no longer a bug. This is the honest σ_geo replacement if a spectral number is ever wanted.
%| 3. **Mean resultant length** (directional statistics) — norm of the average unit vector: ≈1 concentrated, ≈0 spread. Cheapest one-number concentration summary; the uncentered top singular vector alongside it names the shared direction (in the collapse case: the popularity axis).
%|
%| ## The extremely interesting boundary case: at the double tie, norm stops being gauge — and the answer to "which norm?" is ALL of them
%|
%| In the single-sided hosts (fI on the I host, fU on the U host) the gauge argument is exact: prototypes and the prototype-side entity embedding are cosine-only (norms gauge; for composed embeddings the *summand* row norms still matter for direction, but the total norm doesn't); only the linear-side entity embedding has a meaningful norm.
%|
%| At the **double tie** (`prototypes_double_tie`, the fUfI merge target) every vector picks up an inner-product job, so **no norm is pure gauge anymore**:
%| - **Prototype row norms** become meaningful via the second job — each side's prototype matrix is reused as the projection basis of the *other* side's linear half, and in that role prototype row k's norm directly scales coordinate k of the opposite side's linear representation.
%| - **Entity embedding norms** become meaningful because each side's embedding also feeds its linear half (inner-product read), not only the cosine branch.
%|
%| Consequences: (a) the instrument choice is **host-dependent** — spread-on-the-ball for single-sided models, and at the merge a scale-aware spectral measure regains behavioral meaning (σ_geo-like measurements stop being nonsense exactly there); (b) this is the same "second job" mechanism the collapse entry hypothesized as the reason the double-tie stayed healthiest — duplicating a prototype there costs representational rank that accuracy *does* see (coheres with the F-DC01-07 host-specificity insight); (c) the host-dependence itself is a finding-shaped observation for the merge chapter (`fufi-merge`): the merge changes not just the model but *which geometry measurements are meaningful on it*.
%|
%| Scrutiny-evidence-not-thesis applies to anything measured under this entry.
%|
%| [[explanations]] [[evaluation]] [[decisions]] [[paper-understanding]] [[phase-4]]
%|
%| > ⚠ **CORRECTION 2026-08-17:** parts of this entry are corrected in [[2026-08-17_1823_correction-double-tie-projection-basis-and-ui-effective-count]] — the double-tie "prototypes as projection basis" mechanism is code- and paper-refuted (prototype norms stay pure gauge at UI; only entity embeddings hold a second job), and the UI-user "~64 effective" figure is unreproducible (measured: 10 canonical / 8 old run).
% endregion vault: 2026-08-15_1651_spectral-scale-gauge-inversion-direction-diversity-metric
% region vault: 2026-08-17_1823_correction-double-tie-projection-basis-and-ui-effective-count -- auto-copy, edit the vault entry, not this block
%| # CORRECTION: the double-tie "prototypes as projection basis" mechanism is wrong (code- and paper-refuted); July's UI-user effective count is unreproducible
%|
%| Two corrections to the record, found during explanation deep-dive 3.1 (2026-08-17) when Matteo asked where the "duplicating a prototype costs representational rank" sentence came from.
%|
%| ## Correction 1: prototypes do NOT hold a projection job at the double tie
%|
%| The claim (born [[2026-07-15_1112_prototype-collapse-is-a-self-built-popularity-bias]] as "hypothesis-grade", restated as mechanism in [[2026-08-15_1651_spectral-scale-gauge-inversion-direction-diversity-metric]]): at `prototypes_double_tie`, each side's prototype matrix is reused as the projection basis of the other side's linear half, so prototype row norms stop being gauge and duplication costs rank that accuracy sees.
%|
%| **Verified against BOTH sources of truth, 2026-08-17:**
%| - **Code** (`feature_extractor_factories.py` `prototypes_double_tie`, `EmbeddingW`): the projections are free `nn.Linear(d, out_dim)` layers; the ONLY tying is each side's embedding TABLE shared between its two branches (same u feeds cos(u,p_l) and W_u·u).
%| - **Paper** (§3.3): "UI-ProtoMF *defines two linear transformations*, û = W^u·u, t̂ = W^t·t" — free matrices; what is shared "across the two units" is the user/item embeddings.
%|
%| **Corrected statement:** prototype norms are pure gauge in EVERY implemented variant, UI included — prototypes only ever enter through cosine (forward, both regularizers, intrinsic naming). One could unit-normalize all prototype rows in the UI checkpoint and every score/ranking/explanation would be bit-identical. What gains a second job at the double tie is the **entity embedding** (cosine leg + linear leg — its norm IS read by the linear branch); the Aug-15 entry's entity-norm bullet stays correct, its prototype-norm bullet does not. Instrument consequence: row-normalize-first direction metrics are exact (not approximate) for prototypes at UI too. Design-space caveat: a future variant that deliberately ties W := P (e.g. a Route-1-style construction) would create the projection job for real — then, and only then, the scale-aware-spectral discussion applies to prototypes.
%|
%| **Gauge ≠ irrelevant-to-training footnote:** grad of cos w.r.t. p scales ~1/‖p‖, so norms modulate direction-update speed; weight decay keeps them uniform (~0.13 at S0.7), making this roughly symmetric.
%|
%| ## Correction 2: July's "UI-user ≈ 64 effective (healthiest)" is unreproducible
%|
%| The July-15 table lists user_item_proto (user side) ≈ 64 effective directions. Measured 2026-08-17 with the committed instrument (row-normalized, duplicate threshold, cross-checked greedy dedup vs connected components — all agree):
%| - canonical cluster checkpoint (expl_deep_dive 3.1): **10** effective of 76 (clone components 24/9/7/6/6/6/5/5/4/4; 410 pairs >0.999), spectral 4.85, 23 distinct lift names;
%| - older local run (`results/user_item_proto_hm_1_month`, different md5): **8** effective.
%|
%| Neither is near 64 → the July number was a measurement error or was taken on a since-replaced checkpoint. **The qualitative claim survives, weakened:** UI is the healthiest H&M host (10 > 5 > 2 effective; names 23 vs 5 vs 7; UI-item 8/11 with the least-wasteful ratio) — but "≈64" must not be quoted.
%|
%| ## Downgrade: WHY UI is healthier is an OPEN question
%|
%| With the projection-basis mechanism dead, the measured UI health advantage (same near-zero coverage regularizer as the collapsed single hosts — so NOT the regs) has no confirmed mechanism. Leading candidates: (a) the entity-embedding second job — the linear branch forces embeddings to stay linearly discriminative, changing what the herding regularizer finds to herd prototypes into; (b) UI-item side drew small K=11 (+ 5× coverage weight, July's own second factor); (c) unknown. The cosbias2x2 UI arms add evidence; a dedicated probe would be needed for attribution.
%|
%| [[phase-4]] [[explanations]] [[paper-understanding]] [[decisions]]
% endregion vault: 2026-08-17_1823_correction-double-tie-projection-basis-and-ui-effective-count
% region vault: 2026-08-17_1320_popularity-stance-2x2-figure-plan-base-settings-gate -- auto-copy, edit the vault entry, not this block
%| # Popularity stance, 2×2 thesis-figure plan, and the fUfI base-settings gate
%|
%| Follow-up to [[2026-08-17_1128_shifted-cosine-baseline-asymmetry-popularity-channel]] — Matteo's closing thoughts on 1.1 of the explanation deep dive, discussed and refined.
%|
%| ## Stance: popularity in recommendations is NOT a negative thing
%|
%| Bias/popularity signal is legitimate — often it IS the right recommendation. The requirement (Matteo's criterion): popularity must enter through a **disclosed, separable channel** (an explicit item bias b_i is the textbook form: transparent, measurable, removable at will) and must NOT corrupt the embedding space, which this thesis builds to be **mathematically interpretable** (see later thesis sections) and, since it is already transparent, ideally **practically interpretable** too. The failure mode is not "the model uses popularity" — it is "the prototype space becomes semantically meaningless because it is so distorted and bent toward the popularity channel". Primarily an interpretability concern, possibly also a performance concern (to be seen — the 2×2 metrics will say).
%|
%| Operationalization (so the stance is measurable, not aesthetic): corruption indicators = **effective prototype count** (distinct names / clone blocks, e.g. 1.1's "76 prototypes, 5 distinct names") and the **baseline-vs-personalized score share** (e.g. 1.1's 80/20). These are exactly what the cosbias2x2 arms move.
%|
%| ## Thesis figure plan (Popularity Channel section)
%|
%| Grid figure (probably one dataset for the thesis render; both generated):
%| - **Rows: U-Proto, I-Proto only** (2 rows). UI tiles are included only if they DEVIATE from the trend; if they just follow it, UI is reported as numbers in the text, not tiles.
%| - **Columns:** cosine b-off · cosine b-on · shifted-cosine b-off · shifted-cosine b-on
%| - **Tile:** tsne_prototypes.png, **prototype names removed entirely** (irrelevant for this figure).
%| - **Per tile, two reported numbers:**
%|   1. **Effective bias percentage** — the unpersonalized component share, normalized so it counts ONLY rank-relevant mass: U-side c·1ᵀt counts, I-side c·1ᵀu counts as 0 (dead gauge) even though it exists, and b_i counts on BOTH sides when bias is on (so I-side bias-on cells are nonzero — the metric differentiates all 8 cells). Exact normalization (signed components; aggregate over actual top-K recommendations) needs one deliberate spec paragraph at instrument-build time.
%|   2. **HR@10 (test).**
%|   3. **Degeneration number: effective direction count** (row-normalized prototypes, pairwise-cosine duplicate threshold — instrument menu in [[2026-08-15_1651_spectral-scale-gauge-inversion-direction-diversity-metric]]), **flagged clearly as a proxy/approximation, NOT a reliable textbook metric** (threshold-dependent; captures clone-collapse, blind to milder concentration; Matteo directive 2026-08-17). Spectral entropy of the normalized prototype matrix as the threshold-free robustness backup if challenged. **Name-collapse count (distinct lift names) stays OUT of the tiles** — it is an explanation-surface indicator confounded by naming vocabulary/filters (H&M's 9 coarse columns vs ml-1m's tag genome), dataset-specific; report it in prose as the geometry→explanation bridge ("N effective directions → N distinct names"). Note from 1.1 vs 1.2: baseline share and geometry corruption DISSOCIATE (H&M 80%/5 names vs ml-1m 87%/20 names, both n=1 breakdowns) — the two tile numbers are genuinely independent gauges; channel usage is architectural, prototype recruitment is data-dependent (weak taste structure required).
%|
%| Two-tier consequence (Matteo, 2026-08-17): the current 18 cosbias2x2 arms (jobs 7426702–7426722) are the **mechanism tier** — same winning config, knob flipped — clean for geometry/effect claims but UNFAIR for metric comparison (hyperparameters were optimized for the shifted/b-off cell only). **For the actual thesis graphic: run full hyperopt combos at full thesis-level config for every cell** (probably one dataset; incumbent cell included — the current reference is fleet/dev-tier). Queues after model freeze per the standing two-tier numbers policy.
%|
%| Figure-time caveat retained: each t-SNE is its own embedding → within-tile qualitative comparison only (clustered vs spread), one caveat sentence in the caption.
%|
%| ## Commitment: fUfI base-settings gate
%|
%| At the end of the cosbias2x2 experiment: **re-evaluate what the base settings (cosine_type × use_bias) should be for building fUfI-ProtoMF**; if the winner is not the incumbent (shifted, bias-off), queue the necessary reruns. At the end of the thesis's popularity-channel section, **justify exactly that decision** in writing. The 2×2 is thus not only a thesis exhibit — it is the deciding experiment for the merge stage's foundation.
%|
%| [[phase-4]] [[explanations]] [[thesis-writing]] [[experiments]] [[decisions]]
% endregion vault: 2026-08-17_1320_popularity-stance-2x2-figure-plan-base-settings-gate
% region vault: 2026-08-17_1937_base-settings-decision-shifted-bias-off-cosbias2x2-results -- auto-copy, edit the vault entry, not this block
%| # DECISION: the lineage stays on (shifted cosine, bias off) — now evidenced, not inherited. cosbias2x2 results + grid drafts
%|
%| Closes (working-tier) the base-settings question opened in [[2026-08-17_1320_popularity-stance-2x2-figure-plan-base-settings-gate]]. All 18 special-experiment arms landed (jobs 7426702–7426722, mechanism tier: reference winning config, only cosine_type/use_bias flipped, same seed).
%|
%| ## The 2×2 results in brief
%|
%| - **I-side placebo PASSED:** item_proto ref vs cosstd — HR@10 0.4938/0.4947 (hm), 0.5711/0.5709 (ml); effective directions and user-cone statistic identical (0.661 = 0.661). The bias pairs are twins too (0.5807/0.5814, b_i spread 0.368 both). The cosine flip does nothing on the item host, in both bias settings — the loss-invariance theorem, now empirical.
%| - **U-side cosstd: the cone displacement.** Denied the free per-item baseline, the model rebuilt popularity by herding the USER cloud into a cone (mean resultant length 0.07→0.50 hm, 0.11→0.62 ml). Prototypes were NOT freed (hm 5→7 eff, still collapsed; ml 15→11, spectral 5.2→3.2 — worse). Metrics −1 to −3 pts. **The +1 shift is partially protective: it gives popularity a parameter-free outlet; removing it displaces the distortion into the embedding geometry itself.**
%| - **Bias sink is host-dependent, exactly as theory says:** I-host (no covert channel) uses b_i heavily (spread 0.15–0.37), +2.6 pts hm, geometry 2→7 eff. U-host/UI (covert channel available) leave b_i nearly untouched (0.006–0.016) until the shift is ALSO removed, then b_i partially substitutes (0.067–0.12).
%| - **b_u spread = 0.0000 in all nine bias arms** — the user bias never trains under sampled softmax (loss-inert), ninefold direct confirmation of the dead-gauge theory.
%| - **fp-twin nugget:** item_proto-hm bias vs cosstd_bias are loss-equivalent models; metrics twin to the 4th digit (0.5204/0.5205) yet geometry lands at 7 vs 2 effective directions — geometry underdetermined by the metric even WITHIN one config (floating-point lottery). Sibling of [[2026-08-17_1407_geometry-underdetermined-by-selection-metric-rashomon]].
%| - UI is where the shift is most load-bearing: cosstd costs −5 pts (hm 0.659→0.608) — relevant for the fUfI merge host.
%|
%| ## Why we stay on (shifted, bias off)
%|
%| 1. **By Matteo's own criterion** (popularity is fine if disclosed and non-corrupting): cosstd achieves the opposite — popularity relocates INTO the geometry (cone), prototypes stay collapsed, and the cleanly separable, measurable B(t) disclosure line is lost. The shift's channel is the most honest place for popularity to live: visible, separable, measurable to the 4th decimal. Manage and disclose, don't amputate.
%| 2. **Bias-on wins only on the I-host**; neutral-to-negative on U/UI (sink unused while the covert channel exists), and it imports the cold-bias-policy problem (untrained b_i = silent cold-item depressant; cold_eval refuses such configs) into the thesis's central cold story. dc02 requires bias-off; comparability across the lineage survives only with the incumbent.
%| 3. **Archaeology: the incumbent IS the paper's setting, and the cosine half was author-tested.** The authors' original upload has use_bias:0 in every shipped config (base_param, first commit), and contains user_proto_{standard,shifted,shifted_and_div}_hyper_params — an unpublished three-way cosine ablation; their final "chose" configs all settled on shifted. Our cosstd result independently reproduces their unpublished choice, mechanism now explained.
%| 4. Frozen S0 table, all dossiers, both candidates were built on the incumbent — zero comparability cost.
%|
%| **Decision wording:** (shifted, bias-off) confirmed as the lineage standard — evidenced, not inherited. Final ratification with equal-treatment numbers at thesis-section writing, per the standing commitment. What we take instead of a switch: the effective-bias-share instrument as standing disclosure; the coverage regularizer (sim_batch) as the actual promising geometry lever (next candidate arm if any); bias-on documented as the I-host-specific alternative.
%|
%| ## Addendum (same day): the effective-bias instrument is a LOWER BOUND — popularity can still hide inside the prototype structure
%|
%| The annotated grids expose the instrument's blind spot, and the blind spot is itself evidence: the U-host cosstd tile reads **eff. bias 0%** while HR barely moves — popularity did not vanish, it went dark (into the user cone, inside the "personalized" mass where no additive decomposition can see it). The instrument counts only the DISCLOSED additive channels (c·1ᵀt terms + b_i); popularity carried by the geometry itself — the cone, and even in shifted models the clone-amplified cos-part (the recruitment mechanism) — is invisible to it. So:
%|
%| - **eff. bias share = a lower bound on popularity mass**, never "the popularity content";
%| - the geometry indicators (cone statistic, effective direction count) are the complementary detectors for concealed popularity — the two families must always be reported together;
%| - this is the transparency argument self-demonstrated: shifted keeps popularity where the instrument can measure it to the percent; standard conceals it. Caption sentence for the thesis figure.
%|
%| ## Thesis placement (Matteo, 2026-08-17)
%|
%| In the section that PRESENTS ProtoMF, at its end: bring this up — **what is concealed in the shifted cosine** (the structural popularity channel and its side asymmetry) **and how it does not seem to impact much** (the 2×2: metrics robust, placebo on the item side, incumbent not dominated). The two grid sheets carry the visual argument.
%|
%| ## Grid drafts (proto versions, mechanism tier)
%|
%| ![[cosbias2x2_grid_hm_PROTO.png]]
%| ![[cosbias2x2_grid_ml_PROTO.png]]
%|
%| (Also at `Master/temp/cosbias2x2_grid_{hm,ml}_PROTO.png`; generator scripts in session scratchpad.)
%|
%| **Open questions for the final graphics:** (a) chosen dataset still open (hm shows the pathology + mechanisms, ml shows the healthy contrast — possibly both, possibly one + numbers); (b) **fair evaluation is probably needed: fully trained runs, i.e. full hyperopt at thesis-level config per cell** (the current arms are away-game mechanism tier — fine for geometry/effect claims, unfair for metric comparison), per the two-tier consequence already recorded in the figure-plan entry.
%|
%| [[phase-4]] [[decisions]] [[explanations]] [[experiments]] [[thesis-writing]]
% endregion vault: 2026-08-17_1937_base-settings-decision-shifted-bias-off-cosbias2x2-results
% 1721: THE merge-stage mechanism entry — shared rows weld the popularity
% channel to the steering (norm holds two jobs with different masters);
% two-sided refund upgrade; the free-vs-composed x tie/no-tie 2x2 is the
% refund/contamination experiment skeleton (host UI = the free control arm);
% expressivity resolution (noid strictly less; ids equally expressive,
% differently priced, no direction-preserving dilation).
% 1651 + 1823 (correction): which geometry instruments are meaningful is
% HOST-DEPENDENT — at the tie the ENTITY embeddings gain the second job
% (linear leg); prototype norms stay pure gauge in every implemented variant
% (the "prototypes as projection basis" mechanism is refuted — never state it).
% 1320 + 1937: base-settings gate — (shifted, bias-off) confirmed for the
% merge foundation, working-tier; the shift is MOST load-bearing on UI
% (cosstd costs -5 pts) — relevant when choosing the merge host settings;
% re-ratify with equal-treatment numbers at writing time.
% Title provisional (Matteo, 2026-07-14; was "Design"): once the merge
% candidate exists, rename by what the tie structure actually looks like —
% "The Tie" / "The Double Tie" / "The Triple Tie" / however many ties it
% takes to merge our thing.
% [PLANNED] double tie under grounded factors; "check which side is free";
% live intercept decomposes per feature word (LightFM's composed bias re-emerges).
\begin{itemize}
    \item One sentence acknowledging the goofy letter-string names: a \emph{versioning device} — all of these models are really just \emph{Attribute-Aware ProtoMF}. \vault{2026-07-14_1152_outline-refinement-ch2-8-structural-decisions}
% region vault: 2026-07-14_1152_outline-refinement-ch2-8-structural-decisions -- auto-copy, edit the vault entry, not this block
%| # Outline refinement walk (Ch 2–8): merge, granularity principle, delta principle, explanation system placement, afU road
%|
%| Full-day outline-refinement session over the thesis skeleton (`Master/thesis/`); decisions batch-logged. All structural, all reflected in `outline.md` + chapter files; supervisor-relevant forks recorded in the SQ register.
%|
%| 1. **Ch 2+3 merged into "Background and Related Work"** (SQ-7): the old Related-Work middle overlapped the Background teaching material almost 1:1 (feature-aware MF ↔ LightFM section; prototype/concept positioning ↔ x→c→y section), and the genuinely novel neighborhood (prototypes in recommendation) is thin. Merged form = teach-then-position along two axes (interpretability axis / mechanism axis), chapter ends with **The Gap** (pincer close: each axis's failure mode → checkable claim → "the gap is the combination, not an ingredient" → modality-light → hand-off). A later split would be quite hard to untangle — flagged for the supervisor *before* prose.
%| 2. **Granularity principle** (now standing, memory-saved): outline headers follow *content mass*, not chapter parity — few flat sections + content-block bullets; sub-headers added at prose time only if blocks demand dividers. Applied: Data ch → 3 flat sections (datasets = source/contents/why-selected only; ALL decision-shaped material incl. variants, cold splits, canonical feature set, ml-1m genre bags, sparsity positioning lives in Preprocessing & Splits); eval-framework ch → Computational Setup / Comparison Design / Part II placeholder.
%| 3. **The explanation system lives in the eval-framework chapter**, demonstrated on the host ProtoMF (already taught in Background) — the demo deliberately exposes the naming gap empirically (centroid-read prototype names) right before the lineage closes it. **Lineage chapters show DELTA ONLY** ("focus on what is new") — fI: parameter-derived item-prototype names; fU: per-purchase regrouping; no host re-demos.
%| 4. **Delta principle across the lineage:** fI = full-build chapter (what/why/how exactly); fU = delta chapter (the not-a-mirror surprises ARE the content; drift guard on the model section); merge = delta of second order, thinnest by design, don't pad. fI's why-section ("Taste Is Revealed, Not Stated") moved to open the chapter and is annotated as the lineage's justification hub (fU's intro calls back to it).
%| 5. **afU road moved into the fU chapter** (revises the committed lineage outline's step-7 presentation home, `2026-07-11_2351`): the cold-start bell rings in fU's empty-histories section, and the argument (user attributes enter last, weakest, only where nothing else exists) needs only the taste argument, not the merge. Section title announces it: "Users with Empty Histories: The Road to **User-Attribute-Aware Prototypes**" — *user*-attribute because everywhere else "attribute" means item attributes. **afU very likely WILL be built** (Matteo) — written as announced road, not hedge. Merge chapter title now "(a)fUfI-ProtoMF"; its Design section provisionally titled "The Tie" (rename by actual tie count once the candidate exists); one-sentence bullet acknowledging the goofy letter-string names as a pure *versioning device* — all these models are really just Attribute-Aware ProtoMF.
%| 6. **Regularizers demoted from "pure" interpretability penalties** (Matteo): like L1/L2 they plausibly also affect accuracy — an empirical question; **small ablation study planned** (accuracy AND legibility effects). Consistency edit in Part II's motivation.
%| 7. Hidden-effects chapter deliberately left unformed (headers = inventory placeholders; expect regrouping + discarding of unmeasured rungs after experiments).
%|
%| [[thesis-writing]] [[decisions]] [[phase-4]]
% endregion vault: 2026-07-14_1152_outline-refinement-ch2-8-structural-decisions
\end{itemize}

\section{Results}
\label{sec:fufi-results}
% [PLANNED] incl. the length-channel refund test (tax at fI/fU, refund at merge
% — registered hypothesis, hidden-effects ladder rung in the hidden-effects
% chapter). Merged-model cold-user empirics, if any materialize, land HERE
% (the afU conceptual outlook moved to the fU chapter's empty-histories
% section, 2026-07-14).
